% Laget automatisk av elevark/lag_elevark.py fra sim/elektron.html. Endringer her blir overskrevet.
\documentclass[11pt]{article}
\usepackage{elevark}
\begin{document}
\elevarkhode{Fysikk 2 · Spesiell relativitetsteori}{Akselerer et elektron}{https://fysikklab.no/sim/elektron.html}{Simuleringen}

Et elektron i ro akselereres gjennom spenningen $U$ og får den kinetiske energien eU. Hvor fort går det etterpå? Klassisk fysikk og relativitetsteorien gir svært forskjellige svar når spenningen blir stor. Etter akselerasjonen går partiklene gjennom et 3~m langt rør sammen med en lyspuls.

\section*{Oppgaver}

Bruk $m_{\mathrm{e}} = 9{,}11 \cdot 10^{-31}\,\mathrm{kg}$, $e = 1{,}60 \cdot 10^{-19}\,\mathrm{C}$ og $c = 3{,}00 \cdot 10^{8}\,\mathrm{m/s}$.

\begin{oppgaveliste}
\oppgave \fase{gjett}{Gjett}I et røntgenrør akselereres elektroner gjennom 100~kV. Hvor stor andel av lysfarten tror dere elektronene får?
\svar
\oppgave \fase{test}{Regn}Regn ut farten med klassisk fysikk og med relativitetsteorien ved 100~kV. Sjekk i simuleringen.
\svar
\oppgave \fase{test}{Regn}Ved hvilken spenning gir den klassiske formelen at elektronet går med lysfarten? Hva sier relativitetsteorien om farten ved denne spenningen?
\svar
\oppgave \fase{forklar}{Forklar}Bruk $E_{\mathrm{k}} = (γ - 1)mc^{2}$ til å forklare hvorfor ingen partikkel med masse kan nå lysfarten, uansett hvor stor spenningen er.
\svar
\oppgave \fase{gjett}{Gjett}Når er klassisk fysikk «god nok»? Finn spenningen der den klassiske farten er 1~\% for høy. Bytt så til et proton. Blir grensen høyere eller lavere? Hvor mye?
\svar
\oppgave \fase{utfordring}{Utfordring}I LHC ved CERN har protonene den kinetiske energien 6,8~TeV. Regn ut $γ$. Hvor mange meter per sekund er protonene langsommere enn lyset? Tips: Når $γ$ er stor, er $1 - v/c \approx 1/(2γ^{2})$.
\svar
\end{oppgaveliste}

\end{document}
